\chapter{Framework}
\label{ch:framework}

\autoref{ch:analysis} analyzed the sources, patterns, and gaps in enterprise application security data integration. This chapter presents the proposed framework: a reusable blueprint that defines the architecture, data model patterns, connector abstractions, and extension points an organization follows to build its own implementation. The framework is prescriptive but platform-aware; it targets the Databricks lakehouse but separates general patterns from platform-specific choices. \autoref{ch:implementation} demonstrates a concrete instance of this blueprint.

%%%%%%%%%%%
%%% 2.1 %%%
%%%%%%%%%%%

\section{Solution Architecture}
\label{sec:solution-arch}

This section presents the framework's architecture at three levels: technology stack, component structure, and data layer organization. The architecture addresses the requirements identified across the domain analysis: heterogeneous source ingestion, vendor-agnostic normalization, dual \gls{olap}/\gls{oltp} serving, and extensibility for new data sources and analytics.

%%% 2.1.1 %%%

\subsection{Technology Stack}
\label{sec:tech-stack}

The framework targets the Databricks platform on \gls{aws}. A key architectural property of the platform is the separation of storage and compute: data resides in cloud object storage while processing engines scale independently, enabling cost-efficient handling of the bursty, heterogeneous workloads typical of security data integration~\cite{Databricks2026Lakewatch}. This subsection maps platform components to framework roles.

\subsubsection{Delta Lake}

Delta Lake is an open table format that adds \gls{acid} transactions, schema enforcement, and time travel to cloud object storage~\cite{Lee2024}. It provides the storage layer for all three medallion tiers. Schema evolution support allows bronze tables to absorb new source fields without pipeline changes. Time travel enables auditing and rollback, critical for a security data platform where data integrity and reproducibility are non-negotiable. \textcite{Armbrust2021} demonstrated that the lakehouse architecture built on Delta Lake achieves both warehouse-grade governance and data lake flexibility, the combination required by the heterogeneous source landscape analyzed in \autoref{sec:source-integration-summary}.

\subsubsection{Unity Catalog}

Unity Catalog provides unified data governance across the lakehouse~\cite{DatabricksUC2025}. It manages a three-level namespace (\texttt{catalog.schema.table}) that maps directly to the medallion layer organization: one catalog per environment, schemas for each source (bronze) and each domain (silver, gold). Fine-grained access control enforces least-privilege access to security data, which often contains sensitive vulnerability details. Automated lineage tracking records data flow from source through transformations to consumption, supporting auditability requirements. Column-level tags enable classification of sensitive fields such as \gls{cve} descriptions and remediation guidance.

\subsubsection{Lakeflow Declarative Pipelines}

\Gls{dltables} is the pipeline orchestration framework for defining batch and streaming transformations~\cite{DatabricksLDP2025}. Pipelines are declared as Python or \gls{sql} functions with explicit dependencies; the framework resolves execution order automatically. Data quality expectations are embedded directly in pipeline definitions as declarative constraints (e.g., ``severity must be one of critical, high, medium, low''). Records that violate expectations are quarantined without stopping the pipeline, implementing the quarantine pattern described in \autoref{sec:data-integration-patterns}. \Gls{dltables} handles incremental processing natively through change data feed, reducing reprocessing cost for large datasets.

\subsubsection{Lakebase}

Lakebase is a serverless PostgreSQL database that provides the \gls{oltp} serving layer~\cite{DatabricksLakebase2026}. It shares the underlying storage layer with the lakehouse, enabling gold-layer tables to be exposed as \gls{oltp}-queryable tables without data duplication. This architecture satisfies the dual serving requirement identified in \autoref{sec:data-architecture}: \gls{olap} queries run against \gls{sql} warehouses for dashboards and trend analysis, while \gls{oltp} queries run against Lakebase for low-latency operational workloads such as issue tracker integration, remediation state lookups, and serving layer \glspl{api}. Lakebase supports scale-to-zero, reducing cost for bursty workloads. Instant database branching creates isolated copies for development and testing without duplicating data.

\subsubsection{Platform Synergies}

Running all components on a single platform provides governance, lineage, and compute benefits that would require significant integration effort across disparate tools. Unity Catalog governs access uniformly from bronze ingestion through gold serving. Lineage tracks data from source \gls{api} through transformations to dashboard queries. The shared compute and storage model eliminates data movement between analytical and operational stores.

The platform also enables integration with Lakewatch, the Databricks \gls{siem} analyzed in \autoref{sec:lakewatch}. Both systems share Delta Lake storage and Unity Catalog governance. The framework's gold-layer outputs (application risk scores, remediation status) can serve as enrichment signals for Lakewatch threat detection, while Lakewatch's runtime security events could feed back into the framework as an additional finding source. This bidirectional potential reinforces the complementarity identified in the related work analysis.

%%% 2.1.2 %%%

\subsection{Component Design}
\label{sec:component-design}

With the platform components established, the framework organizes into five tiers, illustrated in \autoref{fig:component-design}. Data flows top-to-bottom from sources through the platform to consumers.

\begin{figure}[htbp]
\centering
\begin{tikzpicture}[
    node distance=0.5cm,
    tier/.style={rectangle, draw=black!40, fill=gray!8, minimum width=11cm, minimum height=1.2cm, rounded corners=3pt},
    tierbox/.style={rectangle, draw=black!60, fill=white, minimum width=2.2cm, minimum height=0.7cm, font=\small, inner sep=5pt, rounded corners=3pt},
    sourcebox/.style={rectangle, draw=black!50, fill=white, minimum width=1.6cm, minimum height=0.6cm, font=\footnotesize, rounded corners=2pt},
    extbox/.style={rectangle, draw=black!50, fill=white, minimum width=1.8cm, minimum height=0.6cm, font=\footnotesize, rounded corners=2pt},
    consumer/.style={rectangle, draw=black!50, fill=gray!15, minimum width=1.6cm, minimum height=0.6cm, font=\footnotesize, rounded corners=2pt},
    arrow/.style={-{Stealth[length=2mm]}, thick, black!60},
    tierlabel/.style={font=\small\bfseries, text=black!70},
    bronze/.style={fill={rgb,255:red,159;green,122;blue,52}, fill opacity=0.35, text=black, text opacity=1},
    silver/.style={fill={rgb,255:red,192;green,192;blue,192}, fill opacity=0.4, text=black, text opacity=1},
    gold/.style={fill={rgb,255:red,212;green,175;blue,55}, fill opacity=0.35, text=black, text opacity=1},
]

% === DATA SOURCES TIER ===
\node[tier] (sources-bg) {};
\node[tierlabel, left, align=right] at (sources-bg.west) {Data\\Sources};
\node[sourcebox] at ($(sources-bg.center) + (-3.6cm,0)$) (github) {GitHub};
\node[sourcebox] at ($(sources-bg.center) + (-1.8cm,0)$) (sonar) {SonarQube};
\node[sourcebox] at ($(sources-bg.center) + (0,0)$) (snyk) {Snyk};
\node[sourcebox] at ($(sources-bg.center) + (1.8cm,0)$) (nvd) {NVD};
\node[sourcebox] at ($(sources-bg.center) + (3.6cm,0)$) (other) {...};

% === EXTERNAL INTEGRATION TOOLS TIER ===
\node[tier, below=0.5cm of sources-bg, fill=blue!5] (external-bg) {};
\node[tierlabel, left, align=right] at (external-bg.west) {Integration\\Tools};
\node[extbox] at ($(external-bg.center) + (-3cm,0)$) (fivetran) {Fivetran};
\node[extbox] at ($(external-bg.center) + (-1cm,0)$) (airbyte) {Airbyte};
\node[extbox] at ($(external-bg.center) + (1cm,0)$) (meltano) {Meltano};
\node[extbox] at ($(external-bg.center) + (3cm,0)$) (custom) {...};

% === DATABRICKS PLATFORM ===
% Ingestion Tier
\node[tier, below=1.1cm of external-bg, minimum height=1.5cm] (ingestion-bg) {};
\node[tierlabel, left, align=right] at (ingestion-bg.west) {Ingestion\\Tier};
\node[tierbox, align=center] at ($(ingestion-bg.center) + (-3.5cm,0)$) (partner-conn) {Partner\\Connectors};
\node[tierbox, align=center] at ($(ingestion-bg.center) + (0,0)$) (dlt-conn) {dlt\\Connectors};
\node[tierbox, align=center] at ($(ingestion-bg.center) + (3.5cm,0)$) (custom-api) {Custom\\APIs};

% Processing Tier
\node[tier, below=0.5cm of ingestion-bg, minimum height=1.5cm] (processing-bg) {};
\node[tierlabel, left, align=right] at (processing-bg.west) {Processing\\Tier};
\node[tierbox, bronze, minimum width=2.8cm] at ($(processing-bg.center) + (-3.5cm,0)$) (bronze) {Bronze};
\node[tierbox, silver, minimum width=2.8cm] at (processing-bg.center) (silver) {Silver};
\node[tierbox, gold, minimum width=2.8cm] at ($(processing-bg.center) + (3.5cm,0)$) (gold) {Gold};

% Serving Tier
\node[tier, below=0.5cm of processing-bg, minimum height=1.5cm] (serving-bg) {};
\node[tierlabel, left, align=right] at (serving-bg.west) {Serving\\Tier};
\node[tierbox] at ($(serving-bg.center) + (-1.8cm,0)$) (lakebase) {Lakebase};
\node[tierbox] at ($(serving-bg.center) + (1.8cm,0)$) (restapi) {REST API};

% === DATA CONSUMERS ===
\node[tier, below=0.9cm of serving-bg, fill=gray!15] (consumers-bg) {};
\node[tierlabel, left] at (consumers-bg.west) {Data Consumers};
\node[consumer, fill=white] at ($(consumers-bg.center) + (-2.5cm,0)$) (jira) {Jira};
\node[consumer, fill=white] at ($(consumers-bg.center) + (0,0)$) (aspm) {ASPM};
\node[consumer, fill=white] at ($(consumers-bg.center) + (2.5cm,0)$) (siem) {SIEM};

% === ARROWS ===
% Sources to External Tools
\draw[arrow] (sources-bg.south) -- (external-bg.north);

% External Tools to Ingestion
\draw[arrow] (external-bg.south) -- (ingestion-bg.north);

% Direct bypass arrow (Sources to Ingestion) - goes outside on the right
\draw[arrow] (sources-bg.east) -- ++(0.8cm,0) -- ($(ingestion-bg.east) + (0.8cm,0)$) -- (ingestion-bg.east);

% Ingestion to Processing
\draw[arrow] (ingestion-bg.south) -- (processing-bg.north);

% Processing flow
\draw[arrow] (bronze.east) -- (silver.west);
\draw[arrow] (silver.east) -- (gold.west);

% Processing to Serving
\draw[arrow] (processing-bg.south) -- (serving-bg.north);

% Serving to Consumers
\draw[arrow] (serving-bg.south) -- (consumers-bg.north);

% === DATABRICKS BOUNDARY ===
\begin{scope}[on background layer]
\node[draw=red!30, fill=red!5, rounded corners=5pt, fit=(ingestion-bg)(processing-bg)(serving-bg), inner sep=8pt, label={[font=\footnotesize, text=black]above:Databricks Platform}] {};
\end{scope}

\end{tikzpicture}
\caption{Component Design}
\label{fig:component-design}
\end{figure}


\textbf{Data sources} are the external systems analyzed in \autoref{ch:analysis}: application inventories, \gls{scm} platforms, \gls{cicd} tools, security scanners, and vulnerability databases. They are outside the framework's boundary; the framework consumes their \glspl{api} but does not control them.

\textbf{The ingestion tier} hosts connectors that extract data from sources and land it in the bronze layer. Three connector categories serve different integration scenarios. LakeFlow Connect connectors use the platform's managed ingestion service for declarative, zero-code extraction. \gls{sdk} connectors use source-provided client libraries for programmatic extraction with built-in authentication and pagination. \gls{rest} \gls{api} connectors use the open-source \gls{dltool} library for sources that lack a dedicated \gls{sdk}. All connectors share common concerns: authentication, pagination, rate limiting, and incremental state management. The connector framework in \autoref{sec:connector-framework} defines these patterns.

\textbf{The processing tier} implements the medallion architecture (\autoref{sec:medallion-arch}). Delta Lake provides the storage layer; \gls{dltables} orchestrates the transformations. Bronze stores raw ingested data in source-native schemas. Silver normalizes it into the vendor-agnostic entity and finding model defined in \autoref{sec:data-entities}. Gold computes aggregated metrics, \gls{ml}-enriched scores, and consumption-ready datasets. Unity Catalog governs access and tracks lineage across all three layers.

\textbf{The serving tier} exposes processed data to downstream consumers. \gls{sql} warehouses serve \gls{olap} queries for dashboards and analytics. Lakebase serves \gls{oltp} workloads: low-latency lookups, issue tracker integration, and operational \glspl{api}. A \gls{rest} \gls{api} layer provides \gls{http} access to the \gls{oltp} store.

\textbf{Data consumers} are external systems that read from the serving tier: issue trackers (Jira, ServiceNow), \gls{aspm} dashboards, \gls{siem}/\gls{soar} platforms, and custom reporting tools. Like data sources, consumers are outside the framework boundary.

%%% 2.1.3 %%%

\subsection{Medallion Architecture}
\label{sec:medallion-arch}

The processing tier applies the medallion architecture pattern introduced in \autoref{sec:data-architecture}. Each layer is implemented as a Delta Lake schema within a single Unity Catalog, providing unified governance and cross-layer lineage tracking. \autoref{fig:medallion-design} illustrates the layer structure and the transformations between them.

\begin{figure}[htbp]
\centering
\begin{tikzpicture}[
    node distance=1cm,
    layer/.style={rectangle, draw=black!50, minimum width=11cm, minimum height=2.2cm, rounded corners=3pt},
    schemabox/.style={rectangle, draw=black!40, fill=white, minimum width=2.8cm, minimum height=0.6cm, font=\footnotesize},
    arrow/.style={-{Stealth[length=2.5mm]}, thick, black!60},
    transformation/.style={font=\scriptsize, text=black!70, align=center},
    layerlabel/.style={font=\small\bfseries, text=black!80},
    bronze/.style={fill={rgb,255:red,159;green,122;blue,52}, fill opacity=0.25},
    silver/.style={fill={rgb,255:red,192;green,192;blue,192}, fill opacity=0.3},
    gold/.style={fill={rgb,255:red,212;green,175;blue,55}, fill opacity=0.25},
]

% Bronze Layer
\node[layer, bronze] (bronze-bg) {};
\node[layerlabel] at ($(bronze-bg.north) + (0,-0.3)$) {Bronze Layer};
\node[schemabox] at ($(bronze-bg.center) + (-3.2cm,0)$) {github.*};
\node[schemabox] at ($(bronze-bg.center) + (0,0)$) {servicenow.*};
\node[schemabox] at ($(bronze-bg.center) + (3.2cm,0)$) {sonarqube.*};
\node[font=\scriptsize, text=black!60] at ($(bronze-bg.south) + (0,0.35)$) {Raw data, source schemas preserved};

% Bronze to Silver transformation
\node[transformation, below=0.8cm of bronze-bg] (b2s) {
    \textbf{Normalization}\\
    Schema mapping, type casting\\
    Deduplication, entity resolution
};

% Silver Layer
\node[layer, silver, below=0.8cm of b2s] (silver-bg) {};
\node[layerlabel] at ($(silver-bg.north) + (0,-0.3)$) {Silver Layer};
\node[schemabox] at ($(silver-bg.center) + (-3.2cm,0)$) {Entities};
\node[schemabox] at ($(silver-bg.center) + (0,0)$) {Findings};
\node[schemabox] at ($(silver-bg.center) + (3.2cm,0)$) {Relationships};
\node[font=\scriptsize, text=black!60] at ($(silver-bg.south) + (0,0.35)$) {Unified schema, vendor-agnostic};

% Silver to Gold transformation
\node[transformation, below=0.8cm of silver-bg] (s2g) {
    \textbf{Enrichment \& Aggregation}\\
    ML risk scoring, triage classification\\
    Metric computation, materialization
};

% Gold Layer
\node[layer, gold, below=0.8cm of s2g] (gold-bg) {};
\node[layerlabel] at ($(gold-bg.north) + (0,-0.3)$) {Gold Layer};
\node[schemabox] at ($(gold-bg.center) + (-2.5cm,0)$) {ML Enriched};
\node[schemabox] at ($(gold-bg.center) + (2.5cm,0)$) {Aggregations};
\node[font=\scriptsize, text=black!60] at ($(gold-bg.south) + (0,0.35)$) {Consumption-ready, query-optimized};

% Arrows
\draw[arrow] (bronze-bg.south) -- (b2s.north);
\draw[arrow] (b2s.south) -- (silver-bg.north);
\draw[arrow] (silver-bg.south) -- (s2g.north);
\draw[arrow] (s2g.south) -- (gold-bg.north);

\end{tikzpicture}
\caption{Medallion Layer Design}
\label{fig:medallion-design}
\end{figure}

\subsubsection{Bronze Layer}

The bronze layer stores raw data from each source system with minimal transformation. Each connector writes to its own schema (e.g., \texttt{github}, \texttt{servicenow}, \texttt{sonarqube}), preserving the source's native data structure. Records are appended with standard metadata columns: ingestion timestamp, source system identifier, and batch identifier. No business logic is applied at this layer; the goal is a faithful, auditable copy of source data.

Bronze tables use a schema-on-read approach. New fields from source \gls{api} changes are absorbed through additive schema evolution without breaking existing pipelines. Partitioning by ingestion date enables efficient time-range queries and supports retention policies. Records that fail structural validation at ingestion (malformed \gls{json}, missing required fields) are routed to per-source quarantine tables with diagnostic metadata.

\subsubsection{Silver Layer}

The silver layer is the system of record. Bronze-to-silver transformations normalize heterogeneous source data into the vendor-agnostic domain model defined in \autoref{sec:data-entities}. Three table categories occupy this layer:

\begin{itemize}
    \item \textbf{Entity tables} store normalized dimension data: applications, repositories, teams, commits, pull requests, pipeline runs, dependencies, and branch protection policies. Each entity table uses a surrogate key, a natural key from the source system, and \texttt{valid\_from}/\texttt{valid\_to} timestamps for change tracking.
    \item \textbf{Finding tables} store normalized security findings as fact records. Each finding carries a severity mapped to the canonical scale, a lifecycle status, the source tool identifier, and references to the affected entity (repository, application). Finding tables are organized by category: \gls{sast}, \gls{sca}, secrets, \gls{dast}, containers, and \gls{iac}.
    \item \textbf{Relationship tables} store many-to-many mappings: application-to-repository, finding-to-\gls{cve}, and cross-tool deduplication links.
\end{itemize}

The silver layer applies the data quality patterns from \autoref{sec:data-integration-patterns}: severity harmonization, entity normalization, timestamp standardization, and deduplication. \Gls{dltables} expectations enforce constraints declaratively; records that violate expectations are quarantined rather than propagated.

\subsubsection{Gold Layer}

The gold layer computes consumption-ready datasets from silver data. Two categories of gold tables serve distinct purposes:

\begin{itemize}
    \item \textbf{Aggregation tables} compute metrics at defined grains: application risk scores, team remediation rates, \gls{mttr} by severity, \gls{sla} compliance percentages, and time-series trend roll-ups. These tables power the dashboards and executive reports consumed through the \gls{olap} serving path.
    \item \textbf{\gls{ml}-enriched tables} store model outputs: composite risk scores, false positive predictions, and remediation time estimates. These scores augment the aggregation tables with predictive signals. The \gls{ml} workflow patterns are detailed in \autoref{sec:analytics-patterns}.
\end{itemize}

Gold tables use incremental refresh where possible: new silver records trigger recomputation of only the affected aggregation partitions. Full refresh is reserved for metrics that require global recomputation, such as cross-application percentile rankings.

%%%%%%%%%%%
%%% 2.2 %%%
%%%%%%%%%%%

\section{Data Model}
\label{sec:data-model}

This section maps the conceptual domain model from \autoref{sec:data-entities} to physical table designs across the three medallion layers. It defines reusable schema patterns first, then applies them to produce the concrete entity, finding, and aggregation tables the framework provides.

%%% 2.2.1 %%%

\subsection{Schema Patterns}
\label{sec:schema-patterns}

Schema patterns are reusable templates that standardize how tables are created at each medallion layer. They ensure consistency across the data model and lower the barrier for extending it with new entities or sources. Four patterns cover the framework's needs.

\subsubsection{Bronze Pattern}

Every bronze table follows the same structural convention. The source's native payload is stored alongside a fixed set of metadata columns:

\begin{itemize}
    \item \texttt{\_ingestion\_timestamp} — when the record was ingested (\gls{utc}).
    \item \texttt{\_source\_system} — identifier of the originating tool or platform.
    \item \texttt{\_batch\_id} — unique identifier for the ingestion run, supporting idempotent replays.
    \item \texttt{\_raw\_payload} — the original \gls{json} response or record, preserved for auditability.
\end{itemize}

Additional columns from the source's native schema are included alongside the raw payload through schema-on-read. New fields are absorbed via additive schema evolution. Tables are partitioned by ingestion date to support retention policies and efficient time-range queries.

\subsubsection{Silver Entity Pattern}

Entity tables store normalized dimension data. Each follows a standard structure:

\begin{itemize}
    \item \texttt{id} — surrogate key (auto-generated).
    \item \texttt{natural\_key} — the identifier from the source system (e.g., repository full name, application \texttt{sys\_id}).
    \item \texttt{source\_system} — which tool or platform provided the record.
    \item \texttt{valid\_from} / \texttt{valid\_to} — timestamps for change tracking (Type~2 slowly changing dimension). A \texttt{NULL} value in \texttt{valid\_to} indicates the current version.
    \item Domain-specific columns — attributes defined by the entity type (e.g., \texttt{criticality\_tier} for applications, \texttt{primary\_language} for repositories).
\end{itemize}

The natural key plus source system combination uniquely identifies a record's origin. The surrogate key provides a stable reference for foreign keys even when source identifiers change.

\subsubsection{Silver Finding Pattern}

Finding tables store normalized fact data. Each follows a standard structure:

\begin{itemize}
    \item \texttt{finding\_id} — surrogate key.
    \item \texttt{source\_finding\_id} — the identifier from the source tool.
    \item \texttt{source\_tool} — which scanner produced the finding.
    \item \texttt{repository\_id} — foreign key to the repository entity table.
    \item \texttt{severity} — normalized to the canonical four-level scale (critical, high, medium, low).
    \item \texttt{status} — normalized lifecycle state (open, confirmed, resolved, false\_positive, wontfix).
    \item \texttt{detected\_at} / \texttt{resolved\_at} — \gls{utc} timestamps enabling \gls{mttr} calculation.
    \item \texttt{rule\_id} — the tool-specific rule or check that triggered the finding.
    \item Category-specific columns — attributes unique to the finding type (e.g., \texttt{file\_path} and \texttt{line\_range} for \gls{sast}, \texttt{package\_name} and \texttt{installed\_version} for \gls{sca}).
\end{itemize}

\subsubsection{Relationship Pattern}

Relationship tables store many-to-many mappings between entities. Each contains foreign keys to both sides, a source system indicator, and \texttt{valid\_from}/\texttt{valid\_to} timestamps to track when the relationship was active. No payload columns are included; the table's sole purpose is to link entities. Examples include application-to-repository, finding-to-\gls{cve}, and cross-tool deduplication links.

\subsubsection{Gold Aggregation Pattern}

Aggregation tables follow a grain-metric-period structure:

\begin{itemize}
    \item \textbf{Grain columns} define the level of detail: the entity key (application, team, repository) and time period (day, week, month).
    \item \textbf{Metric columns} store computed values: finding counts by severity, \gls{mttr}, \gls{sla} compliance percentage, risk score.
    \item \textbf{Period columns} store the time window: \texttt{period\_start} and \texttt{period\_end} (\gls{utc}).
    \item \textbf{Refresh metadata}: \texttt{computed\_at} timestamp and refresh strategy indicator (incremental or full).
\end{itemize}

This structure enables consistent querying across all gold tables: filter by grain, aggregate over periods, compare across entities.

%%% 2.2.2 %%%

\subsection{Entity Model}
\label{sec:entity-model}

The silver layer instantiates the schema patterns into concrete tables. This subsection specifies the key tables organized by category: entities (dimensions), findings (facts), reference data, and relationships.

\subsubsection{Entity Tables}

\autoref{tab:entity-tables} lists the entity tables and their key domain-specific columns. All tables include the standard silver entity pattern columns (\texttt{id}, \texttt{natural\_key}, \texttt{source\_system}, \texttt{valid\_from}/\texttt{valid\_to}).

\begin{table}[htbp]
\centering
\caption{Silver Entity Tables}
\label{tab:entity-tables}
\begin{tabularx}{\textwidth}{l l X}
\toprule
\textbf{Table} & \textbf{Source} & \textbf{Key Domain Columns} \\
\midrule
applications & \gls{cmdb} & name, criticality\_tier, lifecycle\_status, business\_unit, compliance\_scope \\
repositories & \gls{scm} & full\_name, primary\_language, visibility, default\_branch, archived, last\_activity\_at \\
teams & \gls{cmdb}/\gls{scm} & name, parent\_team\_id, contact\_email \\
commits & \gls{scm} & sha, author, authored\_at, message, repository\_id \\
pull\_requests & \gls{scm} & number, title, state, author, created\_at, merged\_at, repository\_id \\
pipeline\_runs & \gls{cicd} & pipeline\_id, commit\_sha, status, started\_at, finished\_at, repository\_id \\
dependencies & \gls{sca} & package\_name, installed\_version, license, ecosystem, repository\_id \\
branch\_policies & \gls{scm} & required\_reviewers, require\_status\_checks, enforce\_admins, repository\_id \\
\bottomrule
\end{tabularx}
\end{table}

\subsubsection{Finding Tables}

Finding tables share the silver finding pattern columns and add category-specific attributes. \autoref{tab:finding-tables} summarizes each.

\begin{table}[htbp]
\centering
\caption{Silver Finding Tables}
\label{tab:finding-tables}
\begin{tabularx}{\textwidth}{l X}
\toprule
\textbf{Table} & \textbf{Category-Specific Columns} \\
\midrule
sast\_findings & file\_path, line\_start, line\_end, cwe\_id, language \\
sca\_findings & package\_name, installed\_version, fixed\_version, cve\_id, ecosystem \\
secret\_findings & secret\_type, file\_path, commit\_sha, validity\_status \\
dast\_findings & url, http\_method, parameter, attack\_type \\
container\_findings & image\_name, image\_tag, layer, cve\_id \\
iac\_findings & resource\_type, resource\_name, file\_path, policy\_id, benchmark \\
\bottomrule
\end{tabularx}
\end{table}

Each finding table references its repository through \texttt{repository\_id}. Business application context is derived through the relationship tables rather than stored directly on findings, preserving the many-to-many application-to-repository mapping without duplicating finding records.

\subsubsection{Reference Tables}

Reference tables store external intelligence used for enrichment:

\begin{itemize}
    \item \textbf{vulnerabilities} — \gls{cve} records with description, published date, and \gls{cvss} base score from the \gls{nvd}.
    \item \textbf{epss\_scores} — daily \gls{epss} probabilities per \gls{cve}, enabling time-series analysis of exploitation likelihood.
    \item \textbf{kev\_entries} — \gls{cisa} \gls{kev} catalog entries indicating confirmed active exploitation, with the date added.
\end{itemize}

These tables are refreshed on a scheduled basis from their respective external sources. Findings link to them through \gls{cve} identifiers, forming the three-signal enrichment model described in \autoref{sec:static-appsec}.

\subsubsection{Relationship Tables}

Three relationship tables implement the key mappings identified in the domain model:

\begin{itemize}
    \item \textbf{app\_repo\_mapping} — links applications to repositories (many-to-many). This is the most critical relationship for business context attribution.
    \item \textbf{finding\_cve\_mapping} — links findings to \gls{cve} records, enabling vulnerability enrichment.
    \item \textbf{dedup\_links} — links findings identified as duplicates across tools, preserving traceability to each source report while establishing a canonical finding.
\end{itemize}

%%% 2.2.3 %%%

\subsection{Aggregation Model}
\label{sec:aggregation-model}

The gold layer computes consumption-ready metrics from silver data using the aggregation pattern. Each gold table targets a specific stakeholder need identified in \autoref{sec:data-entities}.

\subsubsection{Application Risk Scores}

The \texttt{app\_risk\_scores} table computes a composite risk metric per application per period. Grain columns are \texttt{application\_id} and time period. Metric columns include: open finding counts by severity, weighted risk score (combining severity, \gls{epss} probability, and application criticality tier), \gls{sla} compliance percentage, and a trend indicator (improving, stable, degrading). This table serves application owners who need a single view of their application's security posture.

\subsubsection{Team Metrics}

The \texttt{team\_metrics} table aggregates remediation performance per team per period. Metrics include: \gls{mttr} by severity, finding closure rate, new vs. resolved finding ratio, and \gls{sla} breach count. Leadership uses these metrics for cross-team comparisons and resource allocation decisions.

\subsubsection{Vulnerability Trends}

The \texttt{vulnerability\_trends} table provides time-series data at configurable grains (daily, weekly, monthly). Metrics include: new findings introduced, findings resolved, net open count, severity distribution shifts, and mean age of open findings. These trends power the longitudinal dashboards that track organizational progress.

\subsubsection{Coverage Analysis}

The \texttt{coverage\_analysis} table identifies gaps in security tool coverage. For each repository, it records which tool categories have produced findings or scan records (indicating active scanning) and which have not. Missing coverage is flagged per category: a repository with no \gls{sast} findings and no \gls{sast} pipeline runs likely lacks static analysis integration. This table enables security teams to prioritize tooling rollout.

\subsubsection{Extension Guide}

Adding a new gold table follows a consistent process: define the grain (which entity, which time period), specify the metrics (what to compute), write a silver-to-gold \gls{dltables} transformation, and configure the refresh strategy (incremental or full). The aggregation pattern ensures all gold tables share a common query interface, so dashboards and reporting tools can consume new tables without structural changes.

%%%%%%%%%%%
%%% 2.3 %%%
%%%%%%%%%%%

\section{Environment and Deployment}
\label{sec:env-deploy}

The solution architecture (\autoref{sec:solution-arch}) and data model (\autoref{sec:data-model}) define the framework's components and data structures. Before these can be populated with connectors or analytics, the deployment environment must be provisioned and the codebase organized. This section covers the foundational setup: deployment strategy, project structure, pipeline orchestration, monitoring, and deployment-level verification. These are one-time decisions that establish the platform; the repeatable patterns for extending the framework with new sources and analytics follow in \autoref{sec:connector-framework} and \autoref{sec:analytics-serving}.

%%% 2.3.1 %%%

\subsection{Deployment Strategy}
\label{sec:deployment-strategy}

The framework uses Databricks Asset Bundles as its deployment mechanism~\cite{DatabricksDAB2025}. A bundle is a declarative configuration that describes Databricks resources (jobs, pipelines, cluster policies, permissions) alongside the source code that implements them. A single \texttt{databricks.yml} configuration file at the project root defines all resource mappings, inter-resource dependencies, and environment-specific parameter overrides. This approach treats the deployment as \gls{iac}: the entire platform configuration is version-controlled, reviewable, and reproducible from a single source.

Three environment targets structure the promotion path: development, staging, and production. Each target overrides a set of parameters: the Unity Catalog catalog name (one catalog per environment, following the namespace design from \autoref{sec:tech-stack}), compute cluster sizing, permission grants, and pipeline scheduling intervals. Development uses smaller clusters and relaxed permissions for rapid iteration. Staging mirrors production configuration for pre-release validation. Production enforces strict access controls, full-scale compute, and scheduled execution.

A \gls{cicd} pipeline automates the build-test-deploy cycle~\cite{DatabricksCICD2025}. On each code change, the pipeline executes three stages: validation checks bundle configuration syntax and resource compatibility, testing runs the project's test suite against the development environment, and deployment promotes the validated bundle to the target environment. Staging and production deployments require an approval gate; only reviewed and tested changes reach shared environments. The pipeline runs on standard \gls{cicd} platforms (GitHub Actions, Azure DevOps, GitLab \gls{cicd}) that invoke the Databricks \gls{cli} for bundle operations.

%%% 2.3.2 %%%

\subsection{Project Structure}
\label{sec:project-structure}

The project follows a monorepo layout: all connectors, transformations, analytics, tests, and deployment configuration reside in a single repository. This simplifies dependency management, enables atomic cross-layer changes, and supports the bundle deployment model described in \autoref{sec:deployment-strategy}. The repository organizes into four top-level directories: \texttt{src/} for pipeline source code (connectors, transformations, analytics), \texttt{tests/} for the test suite, \texttt{config/} for lookup tables and environment-specific settings, and the bundle configuration files at the root.

Within \texttt{src/}, each connector is an independent module containing its own ingestion and transformation logic. Shared utilities (authentication helpers, pagination handlers, normalization functions) reside in a common library that connectors import. Silver-to-gold transformations are grouped by analytic rather than by source, since a single gold table may consume data from multiple connectors. This boundary ensures that adding a new connector requires changes only within its own module and the relevant gold-layer pipelines.

Naming conventions enforce consistency across the platform. Unity Catalog objects follow the pattern \texttt{<env>.<layer>\_<source>.<table\_name>}: for example, \texttt{prod.bronze\_github.code\_scanning\_alerts} or \texttt{prod.silver.sast\_findings}. Pipeline names mirror their source module: \texttt{ingest\_github}, \texttt{transform\_github\_silver}. Column names use \texttt{snake\_case} throughout. These conventions reduce cognitive load and simplify governance policies that rely on name-based access rules.

Configuration management separates secrets from non-sensitive settings. Secrets (\gls{api} tokens, service account credentials) are stored in the platform's secret scope and referenced by name in pipeline code; they never appear in source files or bundle configuration. Non-sensitive settings (severity mapping tables, \gls{sla} thresholds, scheduling intervals) are version-controlled as configuration files and loaded at pipeline runtime. Environment-specific overrides in the bundle configuration select the appropriate catalog, secret scope, and compute target for each deployment.

%%% 2.3.3 %%%

\subsection{Pipeline Orchestration}
\label{sec:orchestration}

Lakeflow Jobs is the orchestration engine that schedules and coordinates pipeline execution~\cite{DatabricksLakeflowJobs2026}. A job groups related tasks into a \gls{dag} with explicit dependencies; the engine executes tasks in dependency order and handles retries on failure. Jobs are defined in the bundle configuration alongside the resources they operate on, making orchestration version-controlled and promotable across environments through the same deployment path described in \autoref{sec:deployment-strategy}.

The framework assigns one job per connector. Each connector job contains two sequential tasks: an ingestion task that extracts data from the source \gls{api} into bronze, and a transformation task that normalizes bronze records into silver. This isolation gives each connector an independent failure domain: a GitHub \gls{api} outage does not block ServiceNow ingestion. Gold-layer analytics run as separate jobs that list their upstream connector jobs as dependencies; the orchestrator starts a gold job only after all required silver data is fresh. \gls{ml} model retraining runs on independent schedules, decoupled from the ingestion cycle.

Source characteristics drive scheduling frequency. High-change sources (\gls{scm} platforms, security scanners producing frequent new findings) run on shorter intervals, typically hourly. Stable sources (application inventory from the \gls{cmdb}) run daily. External enrichment sources follow their respective update cadences: \gls{nvd} and \gls{epss} update daily, \gls{cisa} \gls{kev} updates on weekdays. Gold analytics trigger after their upstream connector jobs complete through explicit job dependencies. Where strict ordering adds unnecessary complexity, such as daily reporting workloads, a schedule offset from the expected ingestion window is an acceptable alternative.

Each task is configured with a retry count and backoff interval for transient failures. If retries exhaust, the task fails and downstream tasks within the same job do not execute. Job-level alerts notify operators of persistent failures through configured channels (email, messaging webhooks). Since connector jobs are independent, a single connector degradation does not cascade. Jobs run in parallel on ephemeral job clusters that are created at job start and terminated at completion, preventing resource contention between connectors and ensuring cost efficiency through right-sized compute per workload.

%%% 2.3.4 %%%

\subsection{Monitoring and Observability}
\label{sec:monitoring}

Orchestration ensures pipelines run on schedule; monitoring ensures they run correctly over time. The platform provides system tables that record job execution history, pipeline events, and compute utilization~\cite{DatabricksObservability2026}. The framework builds on these tables to track three dimensions: data freshness, pipeline health, and data quality trends. All monitoring queries run against the same \gls{sql} warehouse that serves analytics, keeping the operational overhead within the platform.

Data freshness tracks the recency of data at each medallion layer. Each bronze table records the timestamp of its last successful ingestion; each silver table records its last transformation run. A freshness monitor queries these timestamps and flags tables where the age exceeds a configurable threshold. Thresholds reflect the scheduling frequency from \autoref{sec:orchestration}: an hourly connector that has not produced data in two hours is stale; a daily enrichment source that has not refreshed in 36 hours is stale. Freshness violations trigger alerts through the same channels configured for job failures.

Pipeline health aggregates job execution metrics from system tables. Key indicators include: job success rate over a rolling window, mean run duration and its trend (detecting performance degradation before it causes timeouts), task-level retry frequency (high retry rates signal intermittent source issues), and quarantine volume (records routed to quarantine tables per ingestion run). A rising quarantine rate often precedes data quality issues in the silver layer and warrants investigation before downstream aggregations are affected.

Alerting bridges monitoring signals to operational response. The framework configures alerts at two levels. Job-level alerts fire on individual failures or threshold breaches, such as a quarantine rate exceeding a configured percentage of ingested records. System-level alerts fire on aggregate degradation: multiple connectors stale simultaneously, overall pipeline success rate below a threshold, or gold tables not refreshed within their expected window. Alerts are delivered through configured channels (email, messaging webhooks, incident management integrations). The thresholds and channels are defined as configuration, not code, enabling operators to tune sensitivity without redeploying pipelines.

%%% 2.3.5 %%%

\subsection{Testing and Validation}
\label{sec:env-testing}

This subsection covers deployment-level verification: the checks that confirm the environment is correctly provisioned and the codebase meets quality standards. Component-level testing patterns for connectors and analytics are defined separately in \autoref{sec:connector-testing} and \autoref{sec:analytics-testing}.

The \gls{cicd} pipeline enforces quality gates before deployment. Linting validates code style and catches common errors (unused imports, undefined variables, type mismatches). Formatting ensures consistent code layout across contributors. Unit tests verify individual functions (severity mapping lookups, timestamp parsers, schema mapping logic) in isolation from the platform. All gates must pass before the pipeline proceeds to deployment.

After deployment, environment validation confirms that infrastructure prerequisites are in place. Automated checks verify that Unity Catalog objects (catalogs, schemas, tables) exist with expected permissions, compute resources (clusters, \gls{sql} warehouses) are configured and accessible, secret scopes contain the required credentials, and \gls{dltables} pipelines are registered. These checks serve as smoke tests for newly provisioned or updated environments. Failures block pipeline execution until the environment is corrected.

The framework uses pytest markers to link individual tests to requirement identifiers. Each test function carries a marker referencing a specific requirement from \autoref{ch:analysis}, such as \texttt{@pytest.mark.requirement("REQ-SCA-001")}. A traceability matrix is generated automatically from test results, mapping each requirement to its covering tests and their pass/fail status. This pattern applies uniformly across environment, connector (\autoref{sec:connector-testing}), and analytics (\autoref{sec:analytics-testing}) testing. The cross-cutting traceability matrix that summarizes coverage across all components is produced in the reference implementation (\autoref{sec:testing-validation}).

%%%%%%%%%%%
%%% 2.4 %%%
%%%%%%%%%%%

\section{Connector Framework}
\label{sec:connector-framework}

The connector framework defines how data moves from the heterogeneous sources analyzed in \autoref{ch:analysis} into the medallion layers described in \autoref{sec:medallion-arch}. \autoref{sec:component-design} introduced three connector categories (LakeFlow Connect, \gls{sdk}, and \gls{rest} \gls{api} with \gls{dltool}); this section specifies the common abstraction they share, the patterns for landing data in bronze, and the patterns for transforming it to silver. The goal is a repeatable recipe: adding a new source requires implementing a well-defined set of concerns against a new \gls{api}, not redesigning the pipeline.

%%% 2.4.1 %%%

\subsection{Connector Abstraction}
\label{sec:connector-abstraction}

Every connector must handle the same set of cross-cutting concerns, but the three connector categories introduced in \autoref{sec:component-design} differ in how much of this work the developer performs versus the platform or library.

\subsubsection{Connector Categories}

\textbf{LakeFlow Connect connectors} use the platform's managed ingestion service. They are declarative: the connector is configured through the bundle definition, not coded. Authentication, pagination, incremental state, and schema inference are handled by the platform. This category suits sources with supported LakeFlow Connect integrations where full-table ingestion meets requirements. The trade-off is limited flexibility: custom transformation logic, selective field extraction, and non-standard pagination cannot be injected.

\textbf{\gls{sdk} connectors} use a source-provided client library (e.g., the \gls{aws} \gls{sdk} for Python, PyGitHub, python-gitlab) for programmatic extraction. These \glspl{sdk} encapsulate authentication, pagination, rate limiting, and object model mapping, letting the developer work with high-level methods rather than raw \gls{http} requests. This category suits sources that offer an official \gls{sdk} with good coverage of the required endpoints. The developer controls the extraction logic, but the \gls{sdk} handles the transport-level concerns.

\textbf{\gls{rest} \gls{api} connectors with \gls{dltool}} use the open-source \gls{dltool} library to build ingestion pipelines against sources that lack a dedicated \gls{sdk}. The library provides pre-built authentication, pagination, and incremental loading components that the developer composes declaratively against the source's \gls{rest} \gls{api}. This category handles the mechanical concerns (request construction, response parsing, state management) while the developer specifies endpoints, schema mapping, and extraction parameters.

The decision follows a preference order: use LakeFlow Connect if a supported integration exists and meets requirements; use the source's \gls{sdk} when one is available with sufficient endpoint coverage; fall back to \gls{dltool} for sources with only a \gls{rest} \gls{api}. The five responsibilities defined below represent the full contract. \gls{rest} \gls{api} connectors implement them through \gls{dltool} library components. \gls{sdk} connectors delegate them to the client library. LakeFlow Connect connectors satisfy them through platform configuration. The reference implementation in \autoref{ch:implementation} demonstrates all three categories.

\subsubsection{Authentication}

Security tool \glspl{api} use several authentication mechanisms: \glspl{pat}, OAuth~2.0 client credentials, and service account keys. The connector abstraction externalizes credentials through a platform secret scope. Each connector declares which credential type it requires; the framework resolves it at runtime. This separation keeps secrets out of pipeline code and enables credential rotation without redeploying connectors.

\subsubsection{Pagination}

The source integration summary in \autoref{sec:source-integration-summary} shows that \gls{rest} \glspl{api} are the dominant integration surface. These \glspl{api} return results in pages. Two pagination strategies cover the majority of sources: offset-based (page number and size parameters) and cursor-based (an opaque token returned with each response). The connector abstraction encapsulates the pagination strategy so that downstream pipeline logic receives a continuous stream of records regardless of the underlying mechanism. GraphQL-based sources such as the GitHub \gls{api} use cursor-based pagination exclusively.

\subsubsection{Rate Limiting}

Source \glspl{api} enforce rate limits, typically expressed as requests per time window. The connector abstraction implements adaptive backoff: when a rate limit response is received (\gls{http} 429), the connector pauses for the duration specified in the response headers before retrying. For sources without explicit rate limit headers, configurable per-source rate limits prevent exceeding undocumented thresholds. Transient errors (\gls{http} 5xx) trigger exponential backoff with a configurable maximum retry count.

\subsubsection{Incremental State}

Full re-ingestion does not scale for enterprise environments with large data volumes, as discussed in \autoref{sec:data-integration-patterns}. Each connector maintains a high-water mark: the timestamp or cursor of the last successfully ingested record. On subsequent runs, the connector resumes from this position, fetching only new or changed records. The high-water mark is persisted in a state table within the bronze schema, co-located with the data it tracks. For sources that support \gls{cdc} (e.g., ServiceNow's update set mechanism), the connector consumes change events directly.

\subsubsection{Extension Points}

Adding a new source connector follows the category preference order. For a LakeFlow Connect connector: register credentials, add the connector configuration to the bundle definition, and map the landing tables to the bronze schema. For an \gls{sdk} connector: register credentials, implement the extraction logic using the source's client library, define the high-water mark column, and configure the bronze table. For a \gls{rest} \gls{api} connector with \gls{dltool}: register credentials, define the \gls{rest} \gls{api} source configuration (base \gls{url}, endpoints, pagination type), and configure the \gls{dltool} pipeline to land in the bronze schema. In all three cases, no changes to the silver or gold layers are needed at this stage; transformation patterns (\autoref{sec:transformation-patterns}) handle the mapping separately.

%%% 2.4.2 %%%

\subsection{Ingestion Patterns}
\label{sec:ingestion-patterns}

Ingestion moves data from source \glspl{api} into the bronze layer. All connectors follow the same landing pattern, with source-type-specific variations described below.

\subsubsection{Common Landing Pattern}

Every ingestion run produces append-only writes to the target bronze table. The full \gls{api} response is preserved in the \texttt{\_raw\_payload} column alongside the metadata columns defined by the bronze schema pattern (\autoref{sec:schema-patterns}): ingestion timestamp, source system identifier, and batch identifier. No field filtering or transformation occurs at this stage; the bronze layer is a faithful copy of source data.

The batch identifier enables idempotent replays. If a pipeline run fails partway through, re-executing the same batch overwrites only the records from that batch, preventing duplicates. For sources where merge semantics are appropriate (e.g., entity snapshots from a \gls{cmdb}), the connector uses upsert writes keyed on the source record's natural identifier.

Records that fail structural validation at landing, such as malformed \gls{json} or responses with unexpected schema changes, are routed to per-source quarantine tables. Each quarantine record includes the raw payload, the error description, and the batch identifier. This quarantine pattern ensures zero silent data loss: every record retrieved from a source is either successfully landed or isolated with a documented failure reason~\cite{DAMA2017}.

\subsubsection{Source Code Management Sources}

\gls{scm} connectors pull repository metadata, commits, pull requests, and branch protection configurations through \gls{rest} or GraphQL \glspl{api}. These sources produce moderate data volumes with frequent updates. The high-water mark is typically the \texttt{updated\_at} timestamp on the most recently modified record. Webhook-triggered runs complement scheduled pulls for near-real-time ingestion of events such as new pull requests or merged commits.

\subsubsection{Security Scanner Sources}

Security scanner connectors handle the widest format diversity. Server-based tools (SonarQube, Snyk, Dependency-Track) expose \gls{rest} \glspl{api} with paginated \gls{json} responses, similar to \gls{scm} sources. \gls{cli}-based tools (Semgrep, TruffleHog, Trivy) produce output files in \gls{json} or \gls{sarif} format~\cite{OASIS2024} that must be collected from \gls{cicd} pipeline artifacts or uploaded to a staging location. Platform-integrated scanners (GitHub Code Scanning, GitLab Security) expose findings through their host platform's \gls{api}, sharing the same authentication and pagination infrastructure as the \gls{scm} connector for that platform.

For \gls{sarif}-producing tools, the connector preserves the full \gls{sarif} document in the raw payload. \gls{sarif} encodes results, rules, and tool metadata in a single structure, and parsing it into individual findings is deferred to the transformation layer.

\subsubsection{Application Inventory Sources}

\gls{cmdb} connectors differ from other sources in two ways. First, they must traverse relationships: an application record references its owning team, associated repositories, and dependent services through foreign key attributes or relationship tables in the \gls{cmdb}. The connector must resolve these references during ingestion, either by following relationship \glspl{api} or by ingesting related tables separately and joining them in the transformation layer. Second, \gls{cmdb} data includes organization-specific custom attributes (e.g., compliance scope, data classification) that vary across deployments. The schema-on-read approach at the bronze layer absorbs these custom fields without requiring connector changes.

%%% 2.4.3 %%%

\subsection{Transformation Patterns}
\label{sec:transformation-patterns}

Transformations move data from the bronze layer to the silver entity and finding tables defined in \autoref{sec:entity-model}. Each transformation is a \gls{dltables} pipeline that reads from one or more bronze tables and writes to a silver table. Five patterns compose the bronze-to-silver path.

\subsubsection{Schema Mapping}

Each source connector has a corresponding schema mapping that extracts typed columns from the bronze raw payload. The mapping defines which \gls{json} fields correspond to which silver columns, applies type casts (e.g., string timestamps to \gls{utc} datetime), and assigns the surrogate key. Mappings are defined declaratively as column expressions in the \gls{dltables} pipeline, making them easy to review and update when source schemas change.

\subsubsection{Normalization}

Three normalization rules bring heterogeneous source data to a common form:

\begin{itemize}
    \item \textbf{Severity harmonization.} Tools use different severity scales, as described in \autoref{sec:static-appsec}. The framework maps each tool's native scale to the canonical four-level model (critical, high, medium, low) through per-tool lookup tables. These tables are maintained as configuration, not code, enabling updates without pipeline changes.
    \item \textbf{Status normalization.} Finding lifecycle states vary across tools (e.g., SonarQube's ``confirmed'' vs. GitHub's ``open''). A per-tool status mapping translates each tool's states to the canonical five-state model (open, confirmed, resolved, false\_positive, wontfix).
    \item \textbf{Timestamp standardization.} All timestamps are converted to \gls{utc}. Source-specific formats (ISO~8601, Unix epoch, tool-specific strings) are parsed during schema mapping and stored as \gls{utc} datetime columns.
\end{itemize}

\subsubsection{Data Quality Validation}

\Gls{dltables} expectations enforce constraints on every record entering the silver layer. Expectations are declared alongside the transformation logic and checked at runtime. Examples include: severity must be one of the four canonical values, status must be one of the five canonical states, repository foreign keys must reference an existing silver entity, and timestamps must fall within a plausible range. Records that violate expectations are quarantined rather than propagated, consistent with the quarantine pattern applied at ingestion.

\subsubsection{Deduplication}

When multiple tools scan the same repository, overlapping findings must be identified and linked. Three deduplication strategies handle different overlap scenarios:

\begin{itemize}
    \item \textbf{Exact \gls{cve} match.} Two \gls{sca} findings referencing the same \gls{cve} and package in the same repository are duplicates. This is the simplest case.
    \item \textbf{Location-based matching.} Two \gls{sast} findings targeting the same file, line range, and weakness category (mapped via \gls{cwe}) in the same repository are likely duplicates. A confidence threshold accounts for minor line number differences caused by tool-specific parsing.
    \item \textbf{Cross-category correlation.} A secret detected by both a \gls{cli}-based scanner and a platform-integrated scanner at the same file path and commit is a duplicate. Matching uses the secret type and file location.
\end{itemize}

Deduplicated findings are not deleted. Instead, a deduplication link record is created in the \texttt{dedup\_links} relationship table (\autoref{sec:entity-model}), preserving traceability to each source report while establishing a canonical finding for aggregation. Deduplication pairs are enumerated per tool combination in the reference implementation (\autoref{ch:implementation}).

\subsubsection{Business Context Attribution}

The final transformation step links findings to business applications. The \texttt{app\_repo\_mapping} relationship table maps repositories to applications. Since this mapping is many-to-many (one application may span multiple repositories; one repository may serve multiple applications), findings inherit application context through a join rather than a direct foreign key. This join enables the gold-layer aggregations in \autoref{sec:aggregation-model} to compute per-application metrics. The mapping itself is sourced from the \gls{cmdb} connector and can be supplemented by automated repository-to-application inference in the \gls{ml} workflows described in \autoref{sec:analytics-patterns}.

%%% 2.4.4 %%%

\subsection{Testing and Validation}
\label{sec:connector-testing}

Connector tests use the pytest infrastructure and requirement markers established in \autoref{sec:env-testing}. Three additional libraries complete the connector testing stack. An \gls{http} request mocking library (such as \texttt{responses} or \texttt{requests-mock}) intercepts \gls{api} calls to simulate source system behavior without network access. A PySpark DataFrame assertion library (such as \texttt{chispa}) compares transformation outputs against expected results with clear diff reporting. \Gls{dltables} expectations, defined alongside transformation logic, provide runtime data quality validation within pipeline execution. Together, these tools enable isolated, reproducible tests that run in the \gls{cicd} pipeline without depending on live source systems.

Ingestion tests verify the connector abstraction concerns from \autoref{sec:connector-abstraction}. Authentication tests confirm that credentials are resolved from the secret scope and that invalid or expired tokens produce clear error messages rather than silent failures. Pagination tests supply multi-page mock responses and verify that the connector retrieves all pages without data loss or duplication. Rate limit tests simulate \gls{http}~429 responses and verify that the connector pauses and retries according to the backoff policy. Incremental state tests run the connector twice with an intermediate data change and confirm the second run fetches only new records, resuming correctly from the persisted high-water mark. Shared test fixtures provide reusable mock response generators for common \gls{api} patterns.

Transformation tests verify the bronze-to-silver path from \autoref{sec:transformation-patterns}. Schema mapping tests supply known bronze payloads and assert that silver columns have correct types, values, and null handling. Severity normalization tests verify that every value in the tool's native severity scale maps to the correct canonical level; the test dataset must cover all source-specific values, including edge cases such as ``informational'' or ``unknown'' that the tool may produce. Status normalization tests cover all tool-specific lifecycle states similarly. Timestamp tests cover format edge cases: ISO~8601 with and without timezone offsets, Unix epoch in seconds and milliseconds, and tool-specific string formats. Tests are data-driven: each case is a pair of (input bronze record, expected silver record), maintained as \gls{json} fixtures alongside the connector module.

Data quality tests verify \gls{dltables} expectation behavior. For each expectation defined on a silver table, a test supplies a record that violates the constraint and confirms it is routed to the quarantine table. A complementary test supplies a valid record and confirms it passes through. Deduplication tests supply known duplicate pairs (same \gls{cve} and package for \gls{sca}, same file, line range, and \gls{cwe} for \gls{sast}, same file path and commit for secrets) and verify that correct \texttt{dedup\_links} records are created. Negative cases supply similar but distinct findings and verify they are not incorrectly linked.

Every connector must pass tests covering: authentication and credential resolution, pagination across multiple response pages, rate limit backoff, incremental state resume, schema mapping for all ingested endpoints, severity and status normalization for all source-specific values, at least one data quality expectation per silver table, and deduplication for all applicable tool overlap pairs. New connectors are not merged until this suite passes in the \gls{cicd} pipeline.

%%%%%%%%%%%
%%% 2.5 %%%
%%%%%%%%%%%

\section{Analytics and Serving Framework}
\label{sec:analytics-serving}

The data model (\autoref{sec:data-model}) and connector framework (\autoref{sec:connector-framework}) defined how data is structured and ingested. This section addresses the final stage: computing consumption-ready insights from silver data and delivering them to stakeholders. \autoref{sec:aggregation-model} defined \emph{which} gold tables the framework provides; this section defines \emph{how} those computations are structured and how the results reach consumers through analytical, operational, and event-driven channels.

%%% 2.5.1 %%%

\subsection{Analytics Patterns}
\label{sec:analytics-patterns}

Gold-layer computations fall into two categories: rule-based analytics that apply deterministic formulas to silver data, and \gls{ml}-driven analytics that learn patterns from historical data. Both produce outputs that follow the gold aggregation pattern defined in \autoref{sec:schema-patterns}.

\subsubsection{Rule-Based Analytics}

Rule-based analytics implement deterministic business logic as \gls{sql} transformations or \gls{dltables} pipelines. Three computation patterns cover the framework's needs.

\textbf{Composite scoring} combines multiple signals into a single metric. The application risk score, for example, weights open finding counts by severity, multiplies by \gls{epss} exploitation probability, and factors in the application's criticality tier. The formula is configurable: organizations assign weights to each signal based on their risk appetite. \gls{sla} compliance calculations follow the same pattern, comparing finding ages against severity-specific thresholds to produce a compliance percentage.

\textbf{Time-series aggregation} rolls up finding data at daily, weekly, and monthly grains. Each roll-up computes period metrics (new findings, resolved findings, net open count) and derived indicators (period-over-period change, moving averages, severity distribution shifts). These aggregations power the trend dashboards that track organizational progress over time.

\textbf{Threshold classification} assigns categorical labels based on metric values. Risk tiers (critical, high, medium, low) are assigned to applications based on their composite score. Coverage gaps are flagged when a repository lacks scan records for a tool category. \gls{sla} breaches are marked when open finding age exceeds the severity-specific threshold. These classifications enable filtering and alerting in downstream serving.

Two refresh strategies apply to rule-based analytics. \textbf{Incremental refresh} recomputes only the partitions affected by new silver records; this suits most aggregations where each record maps to a single grain partition (e.g., one application, one time period). \textbf{Full refresh} recomputes the entire table; this is necessary for metrics that depend on global state, such as cross-application percentile rankings or organization-wide severity distributions.

\subsubsection{ML-Driven Analytics}

\gls{ml}-driven analytics learn patterns from historical silver and gold data to produce predictive signals that rule-based formulas cannot capture. The framework supports four use cases.

\textbf{Composite risk scoring} predicts which applications carry the highest real risk by combining finding data with development activity signals (commit frequency, dependency age, remediation velocity). Unlike the rule-based risk score that weights predefined signals, the \gls{ml} model learns which combinations of factors historically preceded security incidents.

\textbf{False positive prediction} estimates the probability that a new finding is a false positive, based on patterns in historical triage decisions. Features include the source tool, rule identifier, file type, finding age at triage, and the repository's historical false positive rate for the same rule. High-confidence predictions reduce noise for developers by deprioritizing likely false positives in dashboards and notifications.

\textbf{Remediation time estimation} predicts \gls{mttr} for new findings based on historical resolution data. Features include severity, tool category, repository activity level, team workload, and dependency complexity. These predictions inform \gls{sla} compliance forecasting and resource allocation.

\textbf{Anomaly detection} identifies unusual patterns: unexpected spikes in finding volumes, sudden changes in severity distributions, or new finding categories appearing in previously clean repositories. These signals flag potential tool misconfigurations, new attack surfaces, or integration issues that warrant investigation.

The \gls{ml} workflow uses three platform components. MLflow provides experiment tracking, model versioning, and a model registry that manages the lifecycle from development through staging to production~\cite{MLflow2025}. The Feature Store provides reusable feature tables computed from silver and gold data, ensuring consistent feature definitions across training and inference~\cite{DatabricksFeatureEng2025}. Model Serving exposes registered models as real-time inference endpoints for latency-sensitive predictions.

Model outputs integrate into the gold layer through two paths. Batch predictions (risk scores, false positive probabilities) are written to dedicated columns on gold tables during scheduled pipeline runs, augmenting the rule-based metrics with predictive signals. Real-time predictions (on-demand risk assessments for newly ingested findings) are served through inference endpoints and cached in Lakebase for low-latency access.

\subsubsection{Extension Blueprint}

Adding a new analytics workflow follows a consistent process regardless of type. First, define the business question the analytic answers and identify the stakeholder it serves. Second, determine whether a rule-based formula or an \gls{ml} model is appropriate: rule-based suits well-understood relationships with clear formulas; \gls{ml} suits pattern recognition over complex, high-dimensional data. Third, implement the computation as a \gls{dltables} pipeline (rule-based) or an MLflow-tracked experiment (\gls{ml}). Fourth, wire the output to a gold table following the aggregation pattern, configure the refresh strategy, and register the new table in Unity Catalog for governance and lineage tracking.

%%% 2.5.2 %%%

\subsection{Serving Patterns}
\label{sec:serving-patterns}

The serving layer delivers gold-layer outputs to the data consumers identified in \autoref{sec:component-design}. Three delivery modes address different access patterns and latency requirements.

\subsubsection{Analytical Serving}

Analytical serving targets dashboards, reporting tools, and ad-hoc analysis through the \gls{olap} path. A \gls{sql} warehouse executes queries against gold-layer Delta tables, providing sub-five-second response times for pre-computed aggregations. Materialized views cache frequently accessed query patterns, reducing compute cost for recurring dashboard refreshes.

Stakeholder-specific views organize gold data for each audience: executive risk overviews surface application risk scores and organizational trends, team dashboards present remediation metrics and \gls{sla} compliance, and security engineering views expose tool coverage gaps and deduplication statistics. These views are implemented as \gls{sql} views on top of the gold tables, not separate data copies, ensuring consistency across all consumers.

\subsubsection{Operational Serving}

Operational serving targets low-latency workloads through the \gls{oltp} path. Lakebase exposes gold-layer tables as PostgreSQL-queryable relations with sub-50-millisecond response times~\cite{DatabricksLakebase2026}. Because Lakebase shares the underlying storage layer with the lakehouse, gold tables are accessible without data duplication or synchronization pipelines.

Three operational workloads use this path. Remediation state lookups retrieve the current status, severity, and ownership of individual findings for integration with developer workflows. Operational \glspl{api} expose risk scores, coverage data, and finding details through the PostgREST Data \gls{api}, providing \gls{http} access to any gold table without custom \gls{api} code. \gls{ml} inference results cached in Lakebase serve real-time risk assessments for newly ingested findings. Lakebase's scale-to-zero capability reduces cost during off-peak periods when operational queries are infrequent.

\subsubsection{Event-Driven Serving}

Event-driven serving pushes data to external systems rather than waiting for them to query. This mode closes the loop between data analysis and operational action.

\textbf{Automated issue creation} triggers when new critical findings or \gls{sla} breaches are detected. The framework creates work items in issue trackers (Jira, ServiceNow) through their \glspl{api}, linking back to the finding record in the platform. Idempotency guards prevent duplicate ticket creation on pipeline re-runs.

\textbf{Threshold-based notifications} alert stakeholders when metrics cross configured boundaries: a risk score exceeds a threshold, a new \gls{kev}-listed vulnerability is detected, or \gls{sla} compliance drops below a target percentage. Notifications are delivered through configured channels (email, messaging platforms, webhooks).

\textbf{\gls{siem}/\gls{soar} event feeds} push structured security events to platforms such as Lakewatch (\autoref{sec:lakewatch}). \gls{cdc} on gold tables generates a stream of change events, such as new critical findings, risk score changes, or coverage regressions, that \gls{soar} playbooks consume for automated response. The shared infrastructure with Lakewatch, described in \autoref{sec:tech-stack}, enables this integration without external data movement.

\textbf{Bidirectional synchronization} keeps the framework's view consistent with external systems. Issue tracker status updates (e.g., a Jira ticket moved to ``Done'') flow back into the framework through scheduled polling or webhooks, updating the corresponding finding's lifecycle state in the silver layer. This feedback loop ensures that remediation metrics in the gold layer reflect actual resolution progress.

%%% 2.5.3 %%%

\subsection{Testing and Validation}
\label{sec:analytics-testing}

Analytics and serving tests extend the pytest infrastructure from \autoref{sec:env-testing} with additional components. Gold-layer computation tests use PySpark DataFrame assertions (the same library as connector tests) to validate aggregation outputs against known silver inputs. \Gls{ml} model tests use MLflow's \texttt{evaluate()} \gls{api} to compute standard metrics and \texttt{validate\_evaluation\_results()} to enforce minimum performance thresholds before model promotion~\cite{MLflow2025}. Serving layer tests use \gls{http} client libraries to verify endpoint response correctness and measure latency against target thresholds.

Rule-based analytics tests verify the computations from \autoref{sec:analytics-patterns}. Each test supplies a fixed set of silver records with known values and asserts that the gold table output matches expected metrics exactly. Composite scoring tests cover boundary conditions: zero findings (score should be minimal), all-critical findings (maximum score), mixed severity distributions, and cases where enrichment data (\gls{epss} scores, \gls{kev} status) is missing for some findings. Time-series aggregation tests verify correct period bucketing, period-over-period change calculations, and handling of periods with no activity. Threshold classification tests verify correct tier assignment at each boundary value. Refresh idempotency tests run the same pipeline twice on identical silver input and confirm no duplicate or changed gold records.

\Gls{ml} model tests cover the full lifecycle from \autoref{sec:analytics-patterns}. Training tests verify that models converge on the training dataset and produce metrics (accuracy, precision, recall, F1~score) above configured thresholds. Validation tests use a held-out dataset to check for overfitting; the gap between training and validation metrics must remain within a configured tolerance. Feature freshness tests verify that Feature Store tables used for inference contain data no older than a configured staleness threshold. Prediction drift tests compare recent prediction distributions against a baseline and flag significant statistical shifts. Model promotion tests use MLflow's validation \gls{api} to confirm that a candidate model meets or exceeds the currently registered production model's performance before promotion in the model registry.

Serving layer tests verify the three delivery modes from \autoref{sec:serving-patterns}. Analytical serving tests execute representative dashboard queries against the \gls{sql} warehouse and verify response times fall below the five-second target for pre-computed aggregations. Operational serving tests query Lakebase endpoints and verify sub-50-millisecond response times for single-record lookups and correct \gls{json} response structure from the PostgREST \gls{api}. Event-driven serving tests verify issue creation idempotency (triggering the same finding event twice must not create duplicate tickets), notification delivery for threshold breaches, and bidirectional synchronization consistency (an external status update must propagate to the silver layer within the configured polling interval).

Every gold table must include tests for metric correctness against known inputs, boundary conditions, and refresh idempotency. Every \gls{ml} model must include training convergence, validation accuracy, drift monitoring, and promotion gating tests. Every serving endpoint must include response correctness and latency verification. New analytics or serving configurations are not merged until this suite passes.
